\section{Oorsig}

\subsection{Agtergrond}
Akademia het in 2025 vier koshuise wat aan studente akkommodasie bied. Drie van
hierdie koshuise is in Gerhard straat geleë en een is gevestig op die Leriba
kampus. Al vier koshuise het beperkte kombuisfasiliteite wat studente toelaat om
basiese maaltye voor te berei, maar die meerderheid van studente maak staat op
Spys, die amptelike kosverskaffer van Akademia, vir weeklikse maaltye.

Spys bedien 'n groot aantal studente en personeellede daagliks met maaltye vanaf
die weeklikse spyskaart. Die huidige bestelproses behels dat Spys weekliks 'n
e-pos stuur met die volgende week se spyskaart, waarna studente en personeel per
e-pos hul keuses moet aandui. Hierdie stelsel het verskeie logistieke
uitdagings.

\subsection{Probleemstelling}
Die huidige e-pos gebaseerde bestelproses vir maaltye lei tot
verskeie uitdagings.

\begin{description}
  \item[Bestellingverlies] E-posse raak maklik verlore in inbokse of
    word mis gekyk, wat lei tot onvoltooide of foutiewe bestellings.
  \item[Administratiewe las] Spys personeel moet handmatig elke e-pos
    verwerk om bestellings vas te lê, wat arbeidsintensief is en
    geneigd is tot menslike foute.  \item[Beperkte sigbaarheid] Kliënte kan nie maklik hul
    bestel geskiedenis sien of nagaan watter maaltye vir watter dae
    bestellings geplaas is nie.
  \item[Finansiële kompleksiteit] Die menslike faktor in die
    koppeling van bestellings aan studente- en personeelrekeninge vir
    fakturering is 'n tydrowende proses wat administratiewe foute veroorsaak.
  \item[Beperkte gebruikersvriendelikheid] Die huidige proses laat
    nie toe vir maklike besigtiging, wysigings aan bestellings of
    allergie-aanpassings nie.
\end{description}

\subsection{Doelstelling}
Die doel van die projek is om 'n omvattende digitale oplossing te
skep om die bestuur van die weeklikse spyskaart te vereenvoudig deur
die bogenoemde probleme aan te spreek.

Hierdie projek sal die bestuur van die weeklikse spyskaart en
bestellings vereenvoudig, die administratiewe las verminder, en die
gebruikerservaring verbeter.

Die projek sal uit twee hoofkomponente bestaan:
\begin{itemize}
  \item 'n Mobiele toepassing vir die algemene gebruikers wat etes wil bestel
  \item 'n Web gebaseerde platform vir administratiewe take
\end{itemize}
